\chapter{Projekt}

\section{Wymagania funkcjonalne i niefunkcjonalne}

\textbf{Wymagania funkcjonalne:}
\begin{itemize}
\item Autonomiczna i płynna zmiana jasności świecenia lampy (w 3 zdefiniowanych progach) w zależności od zmian natężenia światła w pomieszczeniu \cite{karta_projektu}.
\item Ciągłe monitorowanie poziomu oświetlenia otoczenia z wykorzystaniem zintegrowanego czujnika natężenia światła (np. fotorezystora) \cite{karta_projektu}.
\item Przetwarzanie i przeliczanie zebranych danych z czujnika za pomocą specjalnego algorytmu wykorzystującego zbiory rozmyte i wnioskowanie (Fuzzy Logic) \cite{karta_projektu}.
\item Prezentacja przetwarzania poprawności danych z czujnika przez algorytm rozmyty z wykorzystaniem wykresów (np. w monitorze portu szeregowego) \cite{karta_projektu}.
\end{itemize}

\textbf{Wymagania niefunkcjonalne:}
\begin{itemize}
\item Oparcie systemu na mikrokontrolerze ESP32 połączonym z modułem oświetleniowym w postaci taśmy LED WS2812 \cite{karta_projektu}.
\item Zapewnienie stabilności działania prototypu oznaczającej pracę bez zawieszania się układu przez co najmniej 1 godzinę ciągłej pracy \cite{karta_projektu}.
\item Eliminacja problemu wizualnego „migotania” lampy na granicy wartości zbiorów rozmytych poprzez zaimplementowanie funkcji uśredniania wyników (np. średnia krocząca z czujnika) oraz właściwe dostrojenie logiki \cite{karta_projektu}.
\item Wykorzystanie w schemacie elektrycznym zewnętrznego zasilacza 5V i kondensatorów ze względu na ryzyko dużego poboru prądu przez taśmę LED \cite{karta_projektu}.
\item Umieszczenie komponentów w bezpiecznej, spasowanej obudowie przygotowanej z wykorzystaniem technologii druku 3D \cite{karta_projektu}.
\item Zmniejszenie zmęczenia oczu użytkownika i ograniczenie zużycia prądu dzięki wprowadzonej optymalizacji \cite{karta_projektu}.
\end{itemize}

\section{Przypadki użycia}


\vspace{0.5cm}
\noindent\textbf{Przypadek użycia:} Automatyczna regulacja jasności\
\textbf{Aktor:} Otoczenie / System\
\textbf{Opis:} System autonomicznie dobiera jasność diod na podstawie zmiany warunków świetlnych w otoczeniu \cite{karta_projektu}.\
\textbf{Warunki wstępne:} Mikrokontroler ESP32 jest podłączony do zasilania, a czujnik natężenia światła zbiera dane.\
\textbf{Przebieg:}
\begin{enumerate}
\item Czujnik wykonuje pomiar jasności otoczenia i przekazuje dane.
\item Dane z czujnika trafiają do algorytmu Fuzzy Logic.
\item Algorytm na podstawie określonych reguł (np. "jeśli jest ciemno, ustaw jasność na maksimum") wylicza sygnał sterujący.
\item Mikrokontroler płynnie dostosowuje natężenie własnego światła na taśmie LED.
\end{enumerate}

\vspace{0.5cm}
\noindent\textbf{Przypadek użycia:} Stabilizacja odczytów (zapobieganie migotaniu)\
\textbf{Aktor:} System\
\textbf{Opis:} Algorytm zabezpiecza lampę przed niepożądanym zjawiskiem migotania światła na granicach zdefiniowanych wartości \cite{karta_projektu}.\
\textbf{Warunki wstępne:} System pracuje w trybie ciągłego odczytywania oświetlenia, a algorytm rozmyty aktywnie działa.\
\textbf{Przebieg:}
\begin{enumerate}
\item Czujnik przekazuje odchylenia spowodowane skrajnymi odczytami światła z zewnątrz.
\item Oprogramowanie wykorzystuje zaimplementowaną funkcję uśredniania wyników (średnią kroczącą z czujnika).
\item System płynnie łagodzi nagłe różnice wartości, poprawiając jakość oświetlenia.
\end{enumerate}

\vspace{0.5cm}
\noindent\textbf{Przypadek użycia:} Monitorowanie przetwarzania danych\
\textbf{Aktor:} Członek zespołu (Twórca)\
\textbf{Opis:} Diagnozowanie poprawności działania danych przetworzonych przez algorytm i czujnik \cite{karta_projektu}.\
\textbf{Warunki wstępne:} Prototyp jest stabilnie podłączony przewodem do komputera wyposażonego w kompatybilne środowisko (np. Arduino IDE, PlatformIO, VS Code).\
\textbf{Przebieg:}
\begin{enumerate}
\item Użytkownik uruchamia monitor portu szeregowego.
\item System strumieniowo przesyła dane pobrane z czujników oraz wyniki wyjściowe wnioskowania algorytmu.
\item Użytkownik obserwuje generowane w czasie rzeczywistym wykresy, upewniając się co do poprawności algorytmu rozmytego.
\end{enumerate}